% GENERATED FILE. Edit canonical lessons, then run npm run generate:books.
% canonical-source-hash: fnv1a64:ddfb5b8d3a7f176f
% canonical-lessons: RU-C436-vozvrat, RU-C436-trebovat, RU-C436-vozmushchatsya, RU-C436-ustranyat, RU-C436-kvitantsiya

\chapter{Getting a refund}
\label{ch:ru-getting-a-refund}
\hlchaptermodality{436}

\begin{chapteropening}
\textbf{By the end of this chapter:} \emph{I can ask for my money back: a refund, demanding, protesting, getting a fault fixed and a receipt.}

Complete the last lesson of chapter 436: I can ask for my money back: a refund, demanding, protesting, getting a fault fixed and a receipt.
\end{chapteropening}

\section[\texorpdfstring{\ru{возврат} (vozvrát)}{возврат (vozvrát)}]{\ru{возврат} (vozvrát) — a refund; a return}

\label{lesson:RU-C436-vozvrat}



\begin{quote}
Before the new one: say the Russian for a tip, then the Russian for to be rude.
\end{quote}

\subsection*{You'll want to know: \ru{возврат}}
\textbf{\ru{возврат}} (\emph{vozvrát}) — \textquotedblleft{}a refund; a return\textquotedblright{}.

\begin{quote}
\ru{Без} \ru{чека} \ru{возврат} \ru{невозможен}.
\end{quote}

\begin{quote}
\emph{Bez chéka vozvrát nevozmózhen.}
\end{quote}

\begin{quote}
Without a receipt a refund is not possible.
\end{quote}

Masculine. \emph{\ru{Возврат} \ru{денег}} — a refund.

\begin{cousinweb}
\emph{\ru{возврат}} is built from \emph{\ru{воз}-} 'back' and \emph{-\ru{врат}} 'a turning', the Church Slavonic shape of native \emph{-\ru{ворот}}, on the root of \emph{\ru{вертеть}} 'to turn' (EtymDB lists \emph{\ru{возврат}}).
\end{cousinweb}

\begin{grammarlens}[title={What you've built}]
One more word for telling, explaining and describing at B1.
\end{grammarlens}

\subsection*{Your turn}
\begin{itemize}
\raggedright
  \item \emph{Say it:} \emph{vozvrát}
  \item \emph{Say it:} \emph{vozvrát}, once more
  \item \emph{Say it:} a sentence of your own with \emph{vozvrát}
  \item \emph{Recall:} say the Russian for a tip, then the Russian for to be rude, then say \emph{vozvrát} again
\end{itemize}

\begin{tcolorbox}[breakable,colback=teal!4,colframe=teal!35!black,title={Before you move on}]
What does \ru{возврат} mean? (\textquotedblleft{}A refund; a return\textquotedblright{}.) Say it once more.
\end{tcolorbox}
\section[\texorpdfstring{\ru{требовать} (trébovat)}{требовать (trébovat)}]{\ru{требовать} (trébovat) — to demand}

\label{lesson:RU-C436-trebovat}



\begin{quote}
Before the new one: say the Russian for to be rude, then the Russian for a refund.
\end{quote}

\subsection*{You'll want to know: \ru{требовать}}
\textbf{\ru{требовать}} (\emph{trébovat}) — \textquotedblleft{}to demand\textquotedblright{}.

\begin{quote}
\ru{Клиент} \ru{требует} \ru{вернуть} \ru{деньги}.
\end{quote}

\begin{quote}
\emph{Kliyént trébuyet vernút déngi.}
\end{quote}

\begin{quote}
The customer demands his money back.
\end{quote}

Imperfective; perfective \emph{\ru{потребовать}}. Present \emph{\ru{требую}}. Takes the genitive or an infinitive.

\begin{cousinweb}
\emph{\ru{требовать}} is usually linked with the old word \emph{\ru{треба}} 'a need; an offering'. EtymDB lists the verb without a source, so the link is hedged; it takes \emph{\ru{треба}} from Church Slavonic \emph{trěba}.
\end{cousinweb}

\begin{grammarlens}[title={What you've built}]
One more word for telling, explaining and describing at B1.
\end{grammarlens}

\subsection*{Your turn}
\begin{itemize}
\raggedright
  \item \emph{Say it:} \emph{trébovat}
  \item \emph{Say it:} \emph{trébovat}, once more
  \item \emph{Say it:} a sentence of your own with \emph{trébovat}
  \item \emph{Recall:} say the Russian for to be rude, then the Russian for a refund, then say \emph{trébovat} again
\end{itemize}

\begin{tcolorbox}[breakable,colback=teal!4,colframe=teal!35!black,title={Before you move on}]
What does \ru{требовать} mean? (\textquotedblleft{}To demand\textquotedblright{}.) Say it once more.
\end{tcolorbox}
\section[\texorpdfstring{\ru{возмущаться} (vozmushchátsya)}{возмущаться (vozmushchátsya)}]{\ru{возмущаться} (vozmushchátsya) — to be indignant, to protest}

\label{lesson:RU-C436-vozmushchatsya}



\begin{quote}
Before the new one: say the Russian for a refund, then the Russian for to demand.
\end{quote}

\subsection*{You'll want to know: \ru{возмущаться}}
\textbf{\ru{возмущаться}} (\emph{vozmushchátsya}) — \textquotedblleft{}to be indignant, to protest\textquotedblright{}.

\begin{quote}
\ru{Пассажиры} \ru{возмущались}, \ru{когда} \ru{рейс} \ru{отменили}.
\end{quote}

\begin{quote}
\emph{Passazhíry vozmushchális, kogdá reys otmeníli.}
\end{quote}

\begin{quote}
The passengers were outraged when the flight was cancelled.
\end{quote}

Imperfective; the perfective is \emph{\ru{возмутиться}}. Takes the instrumental: \emph{\ru{возмущаться} \ru{ценами}}.

\begin{cousinweb}
\emph{\ru{возмущаться}} is built from \emph{\ru{воз}-} 'up', \emph{\ru{мутить}} 'to stir up, to trouble' and \emph{-\ru{ся}}. \emph{\ru{Мутить}} is inherited from Proto-Slavic \emph{mǫtiti} 'to stir, to trouble' (EtymDB).
\end{cousinweb}

\begin{grammarlens}[title={What you've built}]
One more word for telling, explaining and describing at B1.
\end{grammarlens}

\subsection*{Your turn}
\begin{itemize}
\raggedright
  \item \emph{Say it:} \emph{vozmushchátsya}
  \item \emph{Say it:} \emph{vozmushchátsya}, once more
  \item \emph{Say it:} a sentence of your own with \emph{vozmushchátsya}
  \item \emph{Recall:} say the Russian for a refund, then the Russian for to demand, then say \emph{vozmushchátsya} again
\end{itemize}

\begin{tcolorbox}[breakable,colback=teal!4,colframe=teal!35!black,title={Before you move on}]
What does \ru{возмущаться} mean? (\textquotedblleft{}To be indignant, to protest\textquotedblright{}.) Say it once more.
\end{tcolorbox}
\section[\texorpdfstring{\ru{устранять} (ustranyát)}{устранять (ustranyát)}]{\ru{устранять} (ustranyát) — to fix, to eliminate (a fault)}

\label{lesson:RU-C436-ustranyat}



\begin{quote}
Before the new one: say the Russian for to demand, then the Russian for to be indignant.
\end{quote}

\subsection*{You'll want to know: \ru{устранять}}
\textbf{\ru{устранять}} (\emph{ustranyát}) — \textquotedblleft{}to fix, to eliminate (a fault)\textquotedblright{}.

\begin{quote}
\ru{Сантехник} \ru{быстро} \ru{устраняет} \ru{неполадки}.
\end{quote}

\begin{quote}
\emph{Santékhnik býstro ustranyáyet nepoládki.}
\end{quote}

\begin{quote}
The plumber fixes faults quickly.
\end{quote}

Imperfective; perfective \emph{\ru{устранить}}. Formal: \emph{\ru{устранить} \ru{проблему}}.

\begin{cousinweb}
\emph{\ru{устранять}} is built from \emph{\ru{у}-} and \emph{-\ru{странять}}, on \emph{\ru{страна}} in its old sense 'a side': to move something to one side. EtymDB takes \emph{\ru{страна}} from Church Slavonic, from Proto-Slavic \emph{storna} 'side, land'.
\end{cousinweb}

\begin{grammarlens}[title={What you've built}]
One more word for telling, explaining and describing at B1.
\end{grammarlens}

\subsection*{Your turn}
\begin{itemize}
\raggedright
  \item \emph{Say it:} \emph{ustranyát}
  \item \emph{Say it:} \emph{ustranyát}, once more
  \item \emph{Say it:} a sentence of your own with \emph{ustranyát}
  \item \emph{Recall:} say the Russian for to demand, then the Russian for to be indignant, then say \emph{ustranyát} again
\end{itemize}

\begin{tcolorbox}[breakable,colback=teal!4,colframe=teal!35!black,title={Before you move on}]
What does \ru{устранять} mean? (\textquotedblleft{}To fix, to eliminate (a fault)\textquotedblright{}.) Say it once more.
\end{tcolorbox}
\section[\texorpdfstring{\ru{квитанция} (kvitántsiya)}{квитанция (kvitántsiya)}]{\ru{квитанция} (kvitántsiya) — a receipt (for a payment or a service)}

\label{lesson:RU-C436-kvitantsiya}



\begin{quote}
Before the new one: say the Russian for to be indignant, then the Russian for to fix.
\end{quote}

\subsection*{You'll want to know: \ru{квитанция}}
\textbf{\ru{квитанция}} (\emph{kvitántsiya}) — \textquotedblleft{}a receipt (for a payment or a service)\textquotedblright{}.

\begin{quote}
\ru{Кассир} \ru{дал} \ru{мне} \ru{квитанцию}.
\end{quote}

\begin{quote}
\emph{Kassír dal mne kvitántsiyu.}
\end{quote}

\begin{quote}
The cashier gave me a receipt.
\end{quote}

Feminine. A \emph{\ru{квитанция}} is a paper slip for a payment or a service; a till receipt is a \emph{\ru{чек}}.

\begin{cousinweb}
\emph{\ru{квитанция}} was borrowed from Dutch \emph{kwitantie} (EtymDB), from Middle Dutch \emph{quitantie}.
\end{cousinweb}

\begin{grammarlens}[title={What you've built}]
One more word for telling, explaining and describing at B1.
\end{grammarlens}

\subsection*{Your turn}
\begin{itemize}
\raggedright
  \item \emph{Say it:} \emph{kvitántsiya}
  \item \emph{Say it:} \emph{kvitántsiya}, once more
  \item \emph{Say it:} a sentence of your own with \emph{kvitántsiya}
  \item \emph{Recall:} say the Russian for to be indignant, then the Russian for to fix, then say \emph{kvitántsiya} again
\end{itemize}

\begin{tcolorbox}[breakable,colback=teal!4,colframe=teal!35!black,title={Before you move on}]
What does \ru{квитанция} mean? (\textquotedblleft{}A receipt (for a payment or a service)\textquotedblright{}.) Say it once more.
\end{tcolorbox}
